% GENERATED FILE. Edit canonical lessons, then run npm run generate:books.
% canonical-source-hash: fnv1a64:1b9acc8976e2951d
% canonical-lessons: UR-C405-rajistar, UR-C405-shanakht, UR-C405-darban, UR-C405-mangana, UR-C405-bakhshish

\chapter{At the Front Desk}
\label{ch:ur-at-the-front-desk}
\hlchaptermodality{405}

\begin{chapteropening}
\textbf{By the end of this chapter:} \emph{I can check in at a hotel, order something to the room and leave a tip.}

Complete the last lesson of chapter 405: I can check in at a hotel, order something to the room and leave a tip.
\end{chapteropening}

\section[\texorpdfstring{\ur{رجسٹر} (rajisṭar)}{رجسٹر (rajisṭar)}]{\ur{رجسٹر} (rajisṭar) — a register (book)}

\label{lesson:UR-C405-rajistar}



\begin{quote}
Before the new one: say the Urdu for to drip, then the Urdu for a facility.
\end{quote}

\subsection*{You'll want to know: \ur{رجسٹر}}
\textbf{\ur{رجسٹر}} — \emph{rajisṭar} — \textquotedblleft{}a register (book)\textquotedblright{}.

\begin{quote}
\ur{ہوٹل} \ur{والے} \ur{نے} \ur{رجسٹر} \ur{میں} \ur{میرا} \ur{نام} \ur{لکھا۔}
\end{quote}

\begin{quote}
\emph{hoṭal vāle ne rajisṭar meṅ merā nām likhā.}
\end{quote}

\begin{quote}
The man at the hotel wrote my name in the register.
\end{quote}

Masculine.

\begin{cousinweb}
\textbf{\ur{رجسٹر}} (\emph{rajisṭar}) is borrowed from English \emph{register}.
\end{cousinweb}

\begin{grammarlens}[title={What you've built}]
One more word for getting around.
\end{grammarlens}

\subsection*{Your turn}
\begin{itemize}
\raggedright
  \item \emph{Say it:} \emph{rajisṭar}
  \item \emph{Say it:} \emph{rajisṭar}, once more
  \item \emph{Say it:} \emph{hoṭal vāle ne rajisṭar meṅ merā nām likhā.} — The man at the hotel wrote my name in the register.
  \item \emph{Recall:} say the Urdu for to drip, then the Urdu for a facility, then say \emph{rajisṭar} again
\end{itemize}

\begin{tcolorbox}[breakable,colback=teal!4,colframe=teal!35!black,title={Before you move on}]
What does \ur{رجسٹر} mean? (\textquotedblleft{}A register (book)\textquotedblright{}.) Say it once more.
\end{tcolorbox}
\section[\texorpdfstring{\ur{شناخت} (shanākht)}{شناخت (shanākht)}]{\ur{شناخت} (shanākht) — identification; identity}

\label{lesson:UR-C405-shanakht}



\begin{quote}
Before the new one: say the Urdu for a facility, then the Urdu for a register (book).
\end{quote}

\subsection*{You'll want to know: \ur{شناخت}}
\textbf{\ur{شناخت}} — \emph{shanākht} — \textquotedblleft{}identification; identity\textquotedblright{}.

\begin{quote}
\ur{کاؤنٹر} \ur{پر} \ur{شناخت} \ur{کے} \ur{لیے} \ur{پاسپورٹ} \ur{دیجیے۔}
\end{quote}

\begin{quote}
\emph{kāunṭar par shanākht ke liye pāsporṭ dījiye.}
\end{quote}

\begin{quote}
Please show your passport at the desk for identification.
\end{quote}

Feminine.

\begin{cousinweb}
\textbf{\ur{شناخت}} (\emph{shanākht}) is usually traced to Persian \emph{shinākht}, \textquotedblleft{}knowing, recognition\textquotedblright{}.
\end{cousinweb}

\begin{grammarlens}[title={What you've built}]
One more word for getting around.
\end{grammarlens}

\subsection*{Your turn}
\begin{itemize}
\raggedright
  \item \emph{Say it:} \emph{shanākht}
  \item \emph{Say it:} \emph{shanākht}, once more
  \item \emph{Say it:} \emph{kāunṭar par shanākht ke liye pāsporṭ dījiye.} — Please show your passport at the desk for identification.
  \item \emph{Recall:} say the Urdu for a facility, then the Urdu for a register (book), then say \emph{shanākht} again
\end{itemize}

\begin{tcolorbox}[breakable,colback=teal!4,colframe=teal!35!black,title={Before you move on}]
What does \ur{شناخت} mean? (\textquotedblleft{}Identification; identity\textquotedblright{}.) Say it once more.
\end{tcolorbox}
\section[\texorpdfstring{\ur{دربان} (darbān)}{دربان (darbān)}]{\ur{دربان} (darbān) — a doorman}

\label{lesson:UR-C405-darban}



\begin{quote}
Before the new one: say the Urdu for a register (book), then the Urdu for identification.
\end{quote}

\subsection*{You'll want to know: \ur{دربان}}
\textbf{\ur{دربان}} — \emph{darbān} — \textquotedblleft{}a doorman\textquotedblright{}.

\begin{quote}
\ur{دربان} \ur{نے} \ur{ہمارے} \ur{لیے} \ur{دروازہ} \ur{کھولا۔}
\end{quote}

\begin{quote}
\emph{darbān ne hamāre liye darvāzā kholā.}
\end{quote}

\begin{quote}
The doorman opened the door for us.
\end{quote}

Masculine.

\begin{cousinweb}
\textbf{\ur{دربان}} (\emph{darbān}) is built from Persian \emph{dar}, \textquotedblleft{}a door\textquotedblright{}, and \emph{-bān}, \textquotedblleft{}a keeper\textquotedblright{}: it is usually traced to Persian \emph{darbān}, \textquotedblleft{}a doorkeeper\textquotedblright{}.
\end{cousinweb}

\begin{grammarlens}[title={What you've built}]
One more word for getting around.
\end{grammarlens}

\subsection*{Your turn}
\begin{itemize}
\raggedright
  \item \emph{Say it:} \emph{darbān}
  \item \emph{Say it:} \emph{darbān}, once more
  \item \emph{Say it:} \emph{darbān ne hamāre liye darvāzā kholā.} — The doorman opened the door for us.
  \item \emph{Recall:} say the Urdu for a register (book), then the Urdu for identification, then say \emph{darbān} again
\end{itemize}

\begin{tcolorbox}[breakable,colback=teal!4,colframe=teal!35!black,title={Before you move on}]
What does \ur{دربان} mean? (\textquotedblleft{}A doorman\textquotedblright{}.) Say it once more.
\end{tcolorbox}
\section[\texorpdfstring{\ur{منگانا} (mangānā)}{منگانا (mangānā)}]{\ur{منگانا} (mangānā) — to send for; to order (to be brought)}

\label{lesson:UR-C405-mangana}



\begin{quote}
Before the new one: say the Urdu for identification, then the Urdu for a doorman.
\end{quote}

\subsection*{You'll want to know: \ur{منگانا}}
\textbf{\ur{منگانا}} — \emph{mangānā} — \textquotedblleft{}to send for; to order (to be brought)\textquotedblright{}.

\begin{quote}
\ur{ہم} \ur{نے} \ur{کمرے} \ur{میں} \ur{چائے} \ur{منگائی۔}
\end{quote}

\begin{quote}
\emph{ham ne kamre meṅ chāe mangāī.}
\end{quote}

\begin{quote}
We ordered tea to the room.
\end{quote}

Past, with \ur{نے} (he or she …): \ur{اس} \ur{نے} \ur{کھانا} \ur{منگایا}, \emph{us ne khānā mangāyā}.

\begin{cousinweb}
\textbf{\ur{منگانا}} (\emph{mangānā}) is built on \textbf{\ur{مانگنا}}, \textquotedblleft{}to ask for\textquotedblright{}, whose Hindi cousin \emph{māgnā} CDIAL lists under Sanskrit \emph{mārgati}, \textquotedblleft{}seeks\textquotedblright{} (CDIAL 10074): it is to have something brought to you.
\end{cousinweb}

\begin{grammarlens}[title={What you've built}]
One more word for getting around.
\end{grammarlens}

\subsection*{Your turn}
\begin{itemize}
\raggedright
  \item \emph{Say it:} \emph{mangānā}
  \item \emph{Say it:} \emph{mangānā}, once more
  \item \emph{Say it:} \emph{ham ne kamre meṅ chāe mangāī.} — We ordered tea to the room.
  \item \emph{Recall:} say the Urdu for identification, then the Urdu for a doorman, then say \emph{mangānā} again
\end{itemize}

\begin{tcolorbox}[breakable,colback=teal!4,colframe=teal!35!black,title={Before you move on}]
What does \ur{منگانا} mean? (\textquotedblleft{}To send for; to order (to be brought)\textquotedblright{}.) Say it once more.
\end{tcolorbox}
\section[\texorpdfstring{\ur{بخشش} (bakhshish)}{بخشش (bakhshish)}]{\ur{بخشش} (bakhshish) — a tip (for service); a gift}

\label{lesson:UR-C405-bakhshish}



\begin{quote}
Before the new one: say the Urdu for a doorman, then the Urdu for to send for.
\end{quote}

\subsection*{You'll want to know: \ur{بخشش}}
\textbf{\ur{بخشش}} — \emph{bakhshish} — \textquotedblleft{}a tip (for service); a gift\textquotedblright{}.

\begin{quote}
\ur{سارا} \ur{نے} \ur{بیرے} \ur{کو} \ur{بخشش} \ur{دی۔}
\end{quote}

\begin{quote}
\emph{sārā ne bere ko bakhshish dī.}
\end{quote}

\begin{quote}
Sara gave the waiter a tip.
\end{quote}

Feminine.

\begin{cousinweb}
\textbf{\ur{بخشش}} (\emph{bakhshish}) is usually traced to Persian \emph{bakhshish}, \textquotedblleft{}a gift, a giving\textquotedblright{}.
\end{cousinweb}

\begin{grammarlens}[title={What you've built}]
One more word for getting around.
\end{grammarlens}

\subsection*{Your turn}
\begin{itemize}
\raggedright
  \item \emph{Say it:} \emph{bakhshish}
  \item \emph{Say it:} \emph{bakhshish}, once more
  \item \emph{Say it:} \emph{sārā ne bere ko bakhshish dī.} — Sara gave the waiter a tip.
  \item \emph{Recall:} say the Urdu for a doorman, then the Urdu for to send for, then say \emph{bakhshish} again
\end{itemize}

\begin{tcolorbox}[breakable,colback=teal!4,colframe=teal!35!black,title={Before you move on}]
What does \ur{بخشش} mean? (\textquotedblleft{}A tip (for service); a gift\textquotedblright{}.) Say it once more.
\end{tcolorbox}
